\chapter[Modelli selezionati e implementazione]{Modelli selezionati e implementazione sperimentale}
\label{chap:implementazione}

Le famiglie presentate nello stato dell'arte descrivono modi differenti di apprendere la normalità e riconoscere le anomalie. Per confrontarle sul caso di studio occorre tradurre questi principi in pipeline concrete, definendo come vengono elaborate le immagini, quali componenti apprendono dai campioni normali e come viene costruita la mappa di anomalia. Il capitolo sviluppa questo passaggio attraverso i quattro modelli selezionati: \textbf{CCW-SSIM}, \textbf{PaDiM}, \textbf{RD4AD} e \textbf{FastFlow}.

Il punto di partenza comune è costituito dalle immagini di \textbf{MVTec AD}, impiegate per l'adattamento al dominio industriale e per il test. Le rappresentazioni iniziali hanno invece una provenienza diversa: PaDiM, RD4AD e FastFlow utilizzano backbone con pesi pre-addestrati su \textbf{ImageNet}, mentre CCW-SSIM apprende da zero sui campioni normali del materiale. Questa differenza aiuta a comprendere sia le procedure di apprendimento sia i costi delle configurazioni adottate.

\section{Criteri di selezione e ruolo dei modelli}
\label{sec:cap3-criteri}

La scelta dei modelli permette di esaminare quattro famiglie metodologiche discusse nel Capitolo~\ref{chap:stato-arte}:
\begin{itemize}
    \item \textbf{CCW-SSIM} \cite{Berasi_2023} appartiene ai metodi basati sulla ricostruzione. Un autoencoder apprende a riprodurre i campioni normali e la discrepanza tra immagine originale e ricostruita è valutata considerando congiuntamente struttura e colore;
    \item \textbf{PaDiM} \cite{Defard_et_al._2021} è un metodo basato sulla rappresentazione. Le feature estratte da una rete pre-addestrata sono modellate, per ciascuna posizione spaziale, mediante una distribuzione gaussiana multivariata;
    \item \textbf{RD4AD} \cite{Deng_et_al._2022} adotta la \textit{Reverse Distillation}. Un encoder \textit{Teacher} congelato fornisce le rappresentazioni di riferimento, mentre un collo di bottiglia e un decoder \textit{Student} apprendono a ricostruirle sui soli campioni normali;
    \item \textbf{FastFlow} \cite{Yu_et_al._2021} impiega flussi di normalizzazione bidimensionali. Trasformazioni invertibili apprendono la distribuzione delle feature normali preservandone l'organizzazione spaziale.
\end{itemize}

I quattro approcci rappresentano pertanto modalità differenti di costruire un modello della normalità: ricostruzione dell'immagine, stima statistica locale delle feature, ricostruzione distillata delle feature e stima esplicita della loro densità.

L'implementazione adotta lo scenario \textit{one-class} definito nel Capitolo~\ref{chap:introduzione}: la fase di adattamento utilizza esclusivamente immagini prive di difetti. Le immagini normali e anomale della partizione di test sono elaborate nella successiva fase di inferenza; le etichette e le maschere dei difetti sono utilizzate per la valutazione, secondo il protocollo illustrato nel capitolo dedicato ai risultati.

All'interno di questo scenario, l'apprendimento assume forme diverse. CCW-SSIM, RD4AD e FastFlow ottimizzano parametri neurali; PaDiM stima le statistiche gaussiane delle feature normali mantenendo congelato il backbone. Il confronto comprende quindi sia modelli con addestramento iterativo sia un approccio di adattamento statistico. Prima di descriverli singolarmente, è utile chiarire i dati e l'organizzazione sperimentale che condividono.

\section{Ambiente e protocollo implementativo comune}
\label{sec:cap3-ambiente}

\subsection{Dataset MVTec AD e categoria leather}

MVTec AD \cite{Bergmann_et_al._2019} raccoglie immagini industriali ad alta risoluzione, organizzate in categorie di oggetti e texture. Il repository conserva le quindici categorie del dataset; le configurazioni e i manifesti verificabili per i quattro modelli fanno tuttavia riferimento alla categoria \texttt{leather}.

Per tale categoria, la partizione \texttt{train/good} contiene 245 immagini normali. La partizione di test contiene 124 immagini: 32 normali e 92 appartenenti alle tipologie di difetto \texttt{color}, \texttt{cut}, \texttt{fold}, \texttt{glue} e \texttt{poke}. Per le immagini difettose sono disponibili 92 maschere nella directory \texttt{ground\_truth}. La medesima organizzazione dei dati è condivisa da tutte le implementazioni.

\subsection{Organizzazione del repository e pipeline di esecuzione}

Il repository comprende una directory per ciascun modello: \texttt{CCW-SSIM}, \texttt{PaDiM}, \texttt{RD4AD} e \texttt{FastFlow}. Ogni directory contiene un punto di ingresso \texttt{Main.py}, una configurazione dedicata e i moduli per caricamento dei dati, costruzione del modello, adattamento e inferenza. Il notebook \texttt{pipeline\_tesi.ipynb} coordina l'esecuzione.

Il dataset è conservato nella directory centrale \texttt{MVTec\_Data} condivisa tra le singole directory locali dei modelli. L'esecuzione comune è articolata nelle seguenti fasi:
\begin{itemize}
    \item adattamento del modello sui campioni normali presenti in \texttt{train/good};
    \item caricamento del checkpoint o delle statistiche apprese;
    \item inferenza sulle immagini della partizione di test;
    \item salvataggio delle mappe di anomalia e dei metadati di esecuzione.
\end{itemize}

Il modulo \texttt{experiment\_protocol.py} registra manifesti e impronte dei file impiegati, rendendo esplicite la configurazione e la provenienza degli artefatti.

\subsection{Ambiente software e hardware}

Gli ambienti registrati nei manifesti del progetto utilizzano Linux su architettura \texttt{x86\_64} e Python~3.11.15. CCW-SSIM è implementato con TensorFlow~2.15 e le API Keras, mentre PaDiM, RD4AD e FastFlow utilizzano PyTorch~2.13.0 e Torchvision~0.28. Le configurazioni rilevate sono riassunte nella Tabella~\ref{tab:cap3-ambienti}.

\begin{table}[htbp]
    \centering
    \small
    \begin{tabularx}{\textwidth}{|P{2.3cm}|Y|Y|}
        \hline
        \textbf{Modello} & \textbf{Framework} & \textbf{Supporto GPU registrato} \\
        \hline
        CCW-SSIM & TensorFlow 2.15, API Keras & CUDA 12.2, cuDNN 8 \\
        \hline
        PaDiM & PyTorch 2.13.0, Torchvision 0.28 & CUDA 13.0 \\
        \hline
        RD4AD & PyTorch 2.13.0, Torchvision 0.28 & CUDA 13.0 \\
        \hline
        FastFlow & PyTorch 2.13.0, Torchvision 0.28 & CUDA 13.0 \\
        \hline
    \end{tabularx}
    \caption{Ambienti software registrati per le implementazioni.}
    \label{tab:cap3-ambienti}
\end{table}

La gestione dell'ambiente radice è definita dai file \texttt{pyproject.toml} e \texttt{uv.lock}. I manifesti indicano l'impiego di una GPU NVIDIA GeForce RTX~4070~Ti.

\subsection{Gestione di configurazioni, checkpoint e artefatti}

Tutte le pipeline fissano il seme pseudo-casuale a 1024. I parametri sono mantenuti nei rispettivi file di configurazione, mentre i pesi o le statistiche apprese sono salvati in formati specifici per il framework: file HDF5 per CCW-SSIM, file \texttt{pickle} per PaDiM e file PTH per RD4AD e FastFlow.

Le mappe di anomalia sono memorizzate come immagini TIFF in virgola mobile a 32 bit, senza quantizzazione preventiva. A ogni esecuzione sono inoltre associati un manifesto JSON e, per i modelli ottimizzati iterativamente, la cronologia di addestramento in formato CSV. Le mappe costituiscono il collegamento tra implementazione e valutazione: da esse saranno ricavati i punteggi delle immagini e le misure di localizzazione discusse nel Capitolo~\ref{chap:risultati}.

\section{CCW-SSIM}
\label{sec:cap3-ccw}

\subsection{Ricostruzione e confronto percettivo}

CCW-SSIM è un metodo basato sulla ricostruzione \cite{Berasi_2023}. Un autoencoder convoluzionale comprime l'immagine in una rappresentazione latente e la ricostruisce. L'addestramento sui soli campioni normali induce la rete a rappresentare regolarità cromatiche e strutturali proprie del dominio osservato.

Indicando con $x$ l'immagine di ingresso e con $g_{\theta}$ l'autoencoder, la ricostruzione è
\begin{equation}
    \widehat{x} = g_{\theta}(x).
\end{equation}
Per valutare la ricostruzione, il modello combina una misura strutturale, derivata dalla CW-SSIM, con una componente cromatica percettiva. Il confronto richiede quindi una trasformazione dell’immagine che separi luminanza e crominanza.

La trasformazione S-CIELAB converte l'immagine RGB nello spazio XYZ, separa l'informazione in canali opponenti e applica filtri spaziali differenti alle componenti di luminanza e crominanza. Dopo la riconversione in CIELAB si ottengono i canali $L^*$, $a^*$ e $b^*$. Il filtraggio tiene conto della diversa sensibilità del sistema visivo umano alle frequenze spaziali \cite{Berasi_2023}.

La componente strutturale è calcolata su coefficienti complessi ottenuti mediante una piramide steerable. Per due insiemi locali di coefficienti $c_x$ e $c_y$, la CW-SSIM può essere espressa come
\begin{equation}
    S_{\mathrm{CW}}(c_x,c_y) =
    \frac{2\left|\sum_i c_{x,i}c_{y,i}^{*}\right|+K}
    {\sum_i |c_{x,i}|^2+\sum_i |c_{y,i}|^2+K},
    \label{eq:cap3-cwssim}
\end{equation}
dove $c_{y,i}^{*}$ è il complesso coniugato e $K$ è una costante di stabilizzazione. La combinazione con la compatibilità cromatica consente di considerare sia le variazioni di struttura sia quelle di colore.

\subsection{Architettura e configurazione adottata}

L'architettura è definita in \texttt{NNModels.py}. Addestramento e inferenza sono gestiti rispettivamente da \texttt{Training.py} e \texttt{Prediction.py}; il preprocessing percettivo è implementato in \texttt{ColorUtils.py}, \texttt{SpatialFilters.py} e nella directory \texttt{Steerables}. La configurazione seleziona la variante seguente:
\begin{center}
    \texttt{Model\_noise\_skip\_bigger\_old}.
\end{center}

L'encoder utilizza convoluzioni $4\times4$ con canali 32, 64, 128, 256 e 2048 e stride rispettivamente 2, 2, 2, 4 e 4. Le attivazioni sono di tipo LeakyReLU. Il decoder è simmetrico e utilizza convoluzioni trasposte $4\times4$ con canali 256, 128, 64, 32 e 3 e stride 4, 4, 2, 2 e 2. L'ultimo strato ha attivazione lineare. Non sono presenti strati di Batch Normalization o Max-Pooling nella variante configurata.

Questa architettura apprende direttamente dalle immagini normali mediante la perdita \texttt{color\_cwssim\_loss}.\footnote{Il codice comprende anche componenti opzionali basati su VGG, esclusi dalla configurazione adottata.} Il trattamento delle immagini e la formulazione della perdita determinano quindi quali regolarità strutturali e cromatiche guidano l’ottimizzazione.

\subsection{Preprocessing e apprendimento della ricostruzione}

Le immagini sorgente della categoria \texttt{leather} hanno dimensione $1024\times1024$ pixel. Il generatore estrae, da ciascuna immagine normale, 64 patch casuali di dimensione $128\times128$, che formano il batch effettivo. Sono configurate 200 epoche e 245 passi per epoca, corrispondenti al numero di immagini normali disponibili.

Il preprocessing converte RGB in XYZ mediante scikit-image, applica la trasformazione nei tre canali opponenti, filtra ciascun canale con kernel spaziali $5\times5$, riconverte in CIELAB e normalizza i valori dividendo per 102. L'ottimizzazione usa Adam con tasso di apprendimento iniziale $10^{-3}$; uno scheduler moltiplica tale valore per 0,5 ogni 15 epoche. I tensori della pipeline CCW-SSIM sono elaborati in virgola mobile a 64 bit.

La perdita configurata combina la CW-SSIM sul primo canale Lab con l'errore cromatico sui canali $a^*$ e $b^*$. La piramide complessa usa sei scale e sei orientazioni, mentre la finestra locale ha lato 7. I checkpoint sono salvati ogni dieci epoche.

La forma cromatica locale non coincide con la distanza euclidea descritta nella fonte. Indicando con $e_{ab}$ l'errore quadratico medio dei due canali già divisi per 102, il codice di addestramento utilizza
\begin{equation}
 e_{ab}=\frac{(a-\widehat a)^2+(b-\widehat b)^2}{2},\qquad
 C_{\mathrm{train}}=1-\frac{\langle e_{ab}\rangle}{45},\qquad
 \mathcal L=\left\langle 1-S_{\mathrm{CW}}C_{\mathrm{train}}\right\rangle.
\end{equation}
La media cromatica è calcolata sul batch e sulle coordinate spaziali. Nella fonte, invece, la componente è espressa tramite $\Delta E=\sqrt{(a^*-\widehat a^*)^2+(b^*-\widehat b^*)^2}$ e un parametro $k$ \cite{Berasi_2023}. La pipeline locale mantiene la combinazione moltiplicativa di struttura e colore, ma costituisce una variante implementativa e non una traduzione letterale della metrica pubblicata.

\subsection{Costruzione della mappa di anomalia}

L'inferenza procede per porzioni dell'immagine (\textit{patch}). Ogni immagine è suddivisa in patch $256\times256$ con stride di 16 pixel; le ricostruzioni sovrapposte sono ricomposte mediante media. La suddivisione del carico è controllata dal parametro \texttt{BatchSplits}, impostato a 25.

La mappa locale media la CW-SSIM alle altezze 7, 8 e 9 della piramide, su sei orientazioni, e applica alla componente cromatica una media spaziale $7\times7$. Indicando tale media con $\mathcal A_7$, il punteggio salvato è
\begin{equation}
 R=1-\overline S_{\mathrm{CW}}\left(1-10\mathcal A_7(e_{ab})\right).
\end{equation}
Il coefficiente cromatico è quindi diverso da quello dell'addestramento: 10 anziché $1/45$, un rapporto pari a 450. Differiscono anche la modalità di media e le altezze della piramide. Il codice non forza il punteggio nell'intervallo $[0,1]$ prima del salvataggio.

La configurazione applica inoltre l'inversione del primo canale Lab ricostruito mediante \texttt{InvertFirstLabAxis=True}. La sua presenza è accertata, ma non dimostra una necessità del decoder; il test diagnostico discusso nel Capitolo~\ref{chap:risultati} verifica separatamente questo punto. La pipeline registra checkpoint HDF5, cronologia CSV, manifesti JSON e mappe TIFF; la ricostruzione è calcolata in memoria per produrre la mappa. La funzione sperimentale per la predizione dell'intera immagine è presente nel codice, ma non è selezionata dalla configurazione utilizzata.

\section{PaDiM}
\label{sec:cap3-padim}

\subsection{Rappresentazione multiscala e modello della normalità}

PaDiM, acronimo di \textit{Patch Distribution Modeling}, descrive ogni posizione dell'immagine mediante feature estratte a più livelli da una rete convoluzionale pre-addestrata \cite{Defard_et_al._2021}. Le rappresentazioni dei livelli iniziali mantengono maggiore dettaglio spaziale, mentre quelle più profonde codificano informazioni semantiche più astratte. L'allineamento spaziale e la concatenazione producono un vettore $z_{i,j}$ per ogni posizione $(i,j)$.

Il backbone rimane congelato: ImageNet fornisce un'inizializzazione generica dell'estrattore, mentre la distribuzione normale specifica del dominio viene stimata sulle immagini MVTec AD.

Per ogni posizione spaziale, PaDiM stima una gaussiana multivariata
\begin{equation}
    z_{i,j} \sim \mathcal{N}(\mu_{i,j},\Sigma_{i,j}).
\end{equation}
La discrepanza di una feature di test dalla distribuzione normale locale è espressa mediante la distanza di Mahalanobis
\begin{equation}
    d_{i,j} = \sqrt{(z_{i,j}-\mu_{i,j})^{\mathrm{T}}
    \Sigma_{i,j}^{-1}(z_{i,j}-\mu_{i,j})}.
    \label{eq:cap3-padim-mahalanobis}
\end{equation}
La modellazione indipendente delle posizioni conserva l'informazione relativa alla collocazione delle strutture nell'immagine.

\subsection{Estrazione e allineamento delle feature}

Il modello è definito in \texttt{NNModels.py}; adattamento e inferenza sono gestiti da \texttt{Training.py} e \texttt{Prediction.py}. Il backbone configurato è Wide ResNet-50-2 con pesi ImageNet-1K V1 forniti da Torchvision. I parametri del backbone restano congelati per l'intera procedura.

Le immagini sono ridimensionate a $224\times224$ e normalizzate con media e deviazione standard ImageNet. Le feature provengono da \texttt{layer1}, \texttt{layer2} e \texttt{layer3}, con 256, 512 e 1024 canali. Le mappe a risoluzione inferiore non sono semplicemente interpolate: la funzione \texttt{embedding\_concat} usa operazioni \texttt{unfold} e \texttt{fold} per allineare e concatenare le feature in corrispondenza delle patch della mappa a risoluzione maggiore, pari a $56\times56$.

\subsection{Stima delle distribuzioni locali}

Le feature allineate costituiscono la rappresentazione sulla quale viene stimata la normalità. La concatenazione completa produce 1792 canali; per contenere la dimensionalità statistica, l'implementazione seleziona deterministicamente 550 indici di canale mediante \texttt{torch.randperm} con seme 1024. Non viene applicata una proiezione casuale: il sottoinsieme è composto direttamente dai canali selezionati e i relativi indici sono salvati insieme alle statistiche.

Le immagini normali sono elaborate in batch da 32, con due processi di caricamento e ordine non mescolato. La procedura esegue un unico passaggio sul insieme di addestramento, calcola media e covarianza per ogni posizione e aggiunge $0{,}01I$ alla covarianza prima dell'inversione. Non sono presenti ottimizzatore o epoche di aggiornamento dei pesi. Le medie, le inverse delle covarianze e gli indici dei canali sono serializzati in un file \texttt{pickle}.

\subsection{Costruzione della mappa di anomalia}

In inferenza, il preprocessing e il sottoinsieme di canali coincidono con quelli usati durante l'adattamento. Per ogni posizione viene calcolata la distanza riportata nell'Equazione~\ref{eq:cap3-padim-mahalanobis}. La mappa $56\times56$ è riportata alla dimensione originale mediante interpolazione bicubica e regolarizzata con un filtro gaussiano di deviazione standard 4.

La pipeline salva le statistiche in formato \texttt{pickle}, il manifesto JSON e una mappa TIFF in virgola mobile a 32 bit per ciascuna immagine di test.

\section{RD4AD}
\label{sec:cap3-rd4ad}

\subsection{Ricostruzione delle feature mediante distillazione}

Come PaDiM, RD4AD utilizza rappresentazioni pre-addestrate, ma apprende a ricostruirle attraverso una rete neurale. Il modello applica la distillazione della conoscenza in direzione inversa rispetto alla configurazione encoder--encoder tradizionale \cite{Deng_et_al._2022}. Il \textit{Teacher Encoder}, pre-addestrato e congelato, produce rappresentazioni multiscala. Un decoder \textit{Student} riceve una rappresentazione compressa e tenta di ricostruire le feature normali del Teacher procedendo dalle rappresentazioni profonde verso quelle meno profonde.

Durante l'adattamento, la somiglianza tra le rappresentazioni del Teacher e quelle ricostruite dallo Student guida l'ottimizzazione. In inferenza, la discrepanza locale tra le coppie di feature corrispondenti forma la mappa di anomalia.

Tra Teacher e Student è inserito un \textit{One-Class Bottleneck Embedding} (OCBE). Il modulo aggrega le feature provenienti da livelli differenti e restringe il passaggio dell'informazione. L'asimmetria tra encoder congelato, collo di bottiglia addestrabile e decoder riduce la possibilità di trasferire direttamente allo Student tutte le attivazioni del Teacher.

Indicando con $f_k(x)$ la feature del Teacher al livello $k$ e con $\widehat{f}_k(x)$ quella ricostruita, la discrepanza coseno locale impiegata per la localizzazione è
\begin{equation}
    a_k(h,w) = 1-
    \frac{f_k(h,w)^{\mathrm{T}}\widehat{f}_k(h,w)}
    {\lVert f_k(h,w)\rVert_2\lVert\widehat{f}_k(h,w)\rVert_2}.
    \label{eq:cap3-rd-cosine}
\end{equation}

\subsection{Configurazione di Teacher, OCBE e decoder}

L'architettura è definita nel file \texttt{NNModels.py}. L'adattamento è gestito da \texttt{Training.py} e l'inferenza da \texttt{Prediction.py}.

Il Teacher è Wide ResNet-50-2 con pesi ImageNet-1K V1 ed espone le feature di \texttt{layer1}, \texttt{layer2} e \texttt{layer3}. Le immagini sono ridimensionate a $256\times256$ e normalizzate secondo i parametri ImageNet. I pesi del Teacher non vengono aggiornati.

Nel modulo OCBE, la prima feature viene ridotta spazialmente con due convoluzioni $3\times3$ a stride 2, ciascuna seguita da Batch Normalization e ReLU; la seconda è ridotta con un solo blocco analogo. Le due mappe sono concatenate con la terza feature e trasformate da un blocco residuo corrispondente al livello profondo della Wide ResNet. Il decoder utilizza blocchi residui con convoluzioni trasposte. Non sono presenti collegamenti \textit{skip} diretti tra Teacher e Student; le connessioni residue appartengono ai blocchi interni del decoder.

La descrizione riguarda il percorso effettivamente utilizzato per ricostruire le feature.\footnote{La convoluzione $1\times1$ denominata \texttt{conv4} è definita nel codice, ma non viene richiamata dal metodo \texttt{forward} della configurazione corrente.}


\subsection{Apprendimento delle feature normali}

L'addestramento usa batch da 16 immagini, due processi di caricamento e 200 epoche. L'ottimizzatore Adam aggiorna soltanto OCBE e Student Decoder con tasso di apprendimento $5\times10^{-3}$ e coefficienti $\beta_1=0{,}5$, $\beta_2=0{,}999$.

Per ciascuno dei tre livelli, l'implementazione appiattisce tutte le dimensioni non appartenenti al batch, calcola la similarità coseno tra il vettore del Teacher e quello ricostruito e media sul batch. La perdita complessiva è la somma delle tre discrepanze:
\begin{equation}
    \mathcal{L}_{\mathrm{RD}} =
    \sum_{k=1}^{3}\left(1-
    \frac{1}{B}\sum_{b=1}^{B}
    \cos\!\left(\operatorname{vec}(f_k^{(b)}),
    \operatorname{vec}(\widehat{f}_k^{(b)})\right)\right).
\end{equation}

\subsection{Costruzione della mappa di anomalia}

In inferenza, la similarità coseno è invece calcolata lungo i canali per ogni coordinata spaziale, come nell'Equazione~\ref{eq:cap3-rd-cosine}. Le tre mappe sono interpolate bilinearmente alla dimensione originale, sommate e filtrate con una gaussiana di deviazione standard 4.

Il checkpoint PTH contiene i parametri di OCBE e Student Decoder insieme ai metadati necessari; il backbone congelato è ricostruito separatamente dai pesi Torchvision. La pipeline produce inoltre cronologia CSV, manifesto JSON e mappe TIFF a 32 bit.

\section{FastFlow}
\label{sec:cap3-fastflow}

\subsection{Modellazione delle feature con flussi invertibili}

Anche FastFlow opera sulle feature di un estrattore pre-addestrato. La normalità viene descritta apprendendo trasformazioni invertibili: il modello rappresenta la distribuzione delle feature normali mediante flussi di normalizzazione bidimensionali \cite{Yu_et_al._2021}. In luogo di appiattire la rappresentazione, le trasformazioni operano su tensori con dimensioni di canale, altezza e larghezza, conservando la corrispondenza spaziale necessaria alla localizzazione.

Un backbone estrae feature a risoluzioni differenti. Per ciascuna scala, un flusso invertibile trasforma la feature $x$ in una variabile latente $z$ che segue una distribuzione normale standard. La densità nel dominio delle feature può essere calcolata esattamente mediante la formula del cambio di variabile.

Data una trasformazione invertibile $z=f_{\theta}(x)$, vale
\begin{equation}
    \log p_X(x)=\log p_Z(f_{\theta}(x))+
    \log\left|\det\frac{\partial f_{\theta}}{\partial x}\right|.
    \label{eq:cap3-fastflow-cov}
\end{equation}
I blocchi di accoppiamento affine dividono i canali in due parti: una parte condiziona scala e traslazione applicate all'altra. Alternando suddivisioni e permutazioni, tutti i canali partecipano alle trasformazioni pur mantenendo calcolabili l'inversa e il determinante Jacobiano.

\subsection{Estrazione e normalizzazione delle feature}

L'implementazione è un porting locale derivato dal commit Anomalib \texttt{889a67f5}. Non richiede a runtime Lightning, Anomalib, timm o FrEIA: il backbone è fornito da Torchvision e le componenti necessarie dei flussi sono implementate nei moduli del progetto.

La configurazione supporta Wide ResNet-50-2 con pesi ImageNet-1K V1. Il backbone è congelato e produce le feature di \texttt{layer1}, \texttt{layer2} e \texttt{layer3}, aventi rispettivamente 256, 512 e 1024 canali e fattori di riduzione spaziale 4, 8 e 16. Le immagini sono ridimensionate a $256\times256$ e normalizzate con i parametri ImageNet. A ogni scala è applicata una Layer Normalization addestrabile sull'intero tensore $(C,H,W)$.

La variante implementata non include backbone Transformer o Vision Transformer; tali possibilità appartengono alla formulazione generale del metodo, ma non alla configurazione del repository.

\subsection{Configurazione e apprendimento dei flussi}

Per ciascuna delle tre scale è costruito un \texttt{FastFlowBlock} composto da otto passi di accoppiamento invertibile. Poiché \texttt{Conv3x3Only} è disattivato, le sottoreti alternano kernel $3\times3$ e $1\times1$. Il rapporto dei canali nascosti è 1 e il parametro di \textit{clamp} è 2.

Ogni passo esegue un accoppiamento affine, una trasformazione affine globale per canale e una permutazione deterministica inizializzata dal seme. La scala dell'accoppiamento è parametrizzata tramite una tangente iperbolica e un fattore 0,1; la scala globale usa una parametrizzazione \textit{softplus}. Non è presente una fase separata di ActNorm dipendente dai dati.

L'addestramento usa batch da 8 immagini, due processi di caricamento e 200 epoche. Adam aggiorna Layer Normalization e flussi con tasso di apprendimento $10^{-3}$ e decadimento dei pesi $10^{-5}$. Per ogni scala, la funzione obiettivo combina l'energia della variabile latente con il log-determinante Jacobiano:
\begin{equation}
    \mathcal{L}_{\mathrm{FF}}=
    \sum_{k=1}^{3}\frac{1}{B}\sum_{b=1}^{B}
    \left(\frac{1}{2}\lVert z_k^{(b)}\rVert_2^2-J_k^{(b)}\right),
\end{equation}
dove $J_k^{(b)}$ è il log-determinante prodotto dal flusso alla scala $k$. I checkpoint intermedi sono salvati ogni cinque epoche e consentono la ripresa dell'esecuzione.

\subsection{Costruzione della mappa di anomalia}

La predizione utilizza batch da 4 immagini. Per ogni scala, il codice calcola la media sui canali di $-\tfrac{1}{2}z^2$, la trasforma esponenzialmente e ne cambia il segno per ottenere la mappa locale. Le mappe sono interpolate a $256\times256$ e aggregate mediante media; la pipeline le riporta poi alla dimensione originale e applica un filtro gaussiano con deviazione standard 4.

I checkpoint PTH contengono i parametri delle normalizzazioni e dei flussi, mentre il backbone congelato viene ricaricato separatamente. Sono prodotti checkpoint intermedi e finali, cronologia CSV, manifesto JSON e mappe TIFF in virgola mobile a 32 bit.

\section{Sintesi delle configurazioni implementative}
\label{sec:cap3-sintesi}

Le descrizioni dei modelli mostrano come uno stesso insieme di immagini normali possa guidare procedure di apprendimento differenti. La Tabella~\ref{tab:cap3-configurazioni} raccoglie le configurazioni adottate, mettendo in relazione rappresentazione iniziale, dimensione dell’input e adattamento al caso di studio.

\begin{table}[htbp]
    \centering
    \small
    \begin{tabularx}{\textwidth}{|P{1.7cm}|Y|P{2.1cm}|Y|}
        \hline
        \textbf{Modello} & \textbf{Rappresentazione} & \textbf{Input e batch} & \textbf{Adattamento su MVTec AD} \\
        \hline
        CCW-SSIM & Autoencoder addestrato da zero & patch $128\times128$, batch 64 & 200 epoche, sole immagini normali \\
        \hline
        PaDiM & Wide ResNet-50-2 ImageNet congelata & $224\times224$, batch 32 & un passaggio per la stima statistica \\
        \hline
        RD4AD & Teacher ImageNet congelato, OCBE e decoder addestrabili & $256\times256$, batch 16 & 200 epoche, sole immagini normali \\
        \hline
        FastFlow & Wide ResNet-50-2 ImageNet congelata e flussi addestrabili & $256\times256$, batch 8 & 200 epoche, sole immagini normali \\
        \hline
    \end{tabularx}
    \caption{Sintesi delle configurazioni implementative presenti nel repository.}
    \label{tab:cap3-configurazioni}
\end{table}

Le differenze di risoluzione, rappresentazione e ottimizzazione sono parte delle condizioni del confronto. CCW-SSIM apprende la ricostruzione dell’immagine, PaDiM stima distribuzioni locali, RD4AD ricostruisce le feature del Teacher e FastFlow ne apprende trasformazioni invertibili. Tutte le pipeline producono infine mappe riferite alle immagini del medesimo insieme di test, sulle quali viene applicato il protocollo comune di valutazione.

Il capitolo seguente esamina quanto queste mappe permettano di distinguere immagini normali e difettose e di localizzare le anomalie. Il confronto quantitativo sarà affiancato dalla lettura degli esempi visivi, per mettere in relazione le prestazioni complessive con risposte incomplete, diffuse o presenti nelle regioni normali. Le differenze implementative qui descritte forniranno il contesto per interpretare gli esiti, mantenendo distinte le osservazioni dalle ipotesi sulle loro cause.
